\documentclass[10pt,twocolumn]{article}
\usepackage[scaled]{helvet}
\renewcommand\familydefault{\sfdefault}
\usepackage[T1]{fontenc}
\usepackage[margin=0.7in]{geometry}
\usepackage{titlesec}
\usepackage{enumitem}
\usepackage{parskip}
\setlength{\columnsep}{0.24in}
\setlist[itemize]{leftmargin=1.2em, topsep=0.2em, itemsep=0.18em}
\titleformat{\section}{\large\bfseries}{\thesection}{1em}{}

\begin{document}
\title{\vspace{-1.6cm}AWS CodePipeline}
\author{AWS ML Study Notes}
\date{}
\maketitle
\small

\section*{Overview}
CodePipeline is AWS's managed continuous delivery orchestrator. It connects source, build, test, approval, and deployment actions into a stage based release flow.

Its main value is disciplined automation. Instead of relying on a person to remember deployment steps, the pipeline executes the same controlled sequence every time.

\section*{Core Concepts}
\begin{itemize}
\item A pipeline is composed of stages, and each stage contains actions such as source retrieval, CodeBuild execution, manual approval, or deployment integration.
\item Artifacts move from one stage to the next so that the exact built output is what gets tested and deployed.
\item Pipelines integrate well with IAM, KMS, EventBridge, and downstream deployment targets such as ECS, Lambda, CloudFormation, or custom actions.
\end{itemize}

\section*{How It Works}
\begin{itemize}
\item A code or configuration change triggers the pipeline, which pulls the source and passes it into automated build and validation steps.
\item Successful outputs are packaged as artifacts, optionally gated by manual approval, then deployed to the target environment through a controlled action.
\item Notifications and audit trails make it clear which version moved where and whether a failure happened in source, build, test, or deployment.
\end{itemize}

\section*{AI and ML Relevance}
\begin{itemize}
\item In ML systems, CodePipeline is often used for infrastructure as code, deployment of inference services, packaging of training images, or promotion of approved model packages.
\item It is particularly useful when the broader platform standardizes on AWS native CI and CD rather than an external orchestration tool.
\end{itemize}

\section*{Operational Notes}
\begin{itemize}
\item Keep pipeline stages purposeful. Mixing unrelated validation and deployment logic into one opaque step makes failures harder to debug.
\item Store artifacts securely and make the pipeline idempotent so reruns do not create confusing partial state.
\item Use manual approval only where human judgment adds real value. Too many approval steps turn automation into ceremony.
\end{itemize}

\section*{What To Remember}
\begin{itemize}
\item CodePipeline is a release conveyor belt. It is most valuable when release steps are standardized and auditable.
\item Artifacts matter because they preserve what was actually built and tested.
\item It solves software promotion, not experiment tracking. That is an important boundary in ML systems.
\end{itemize}

\end{document}
